%---------------------------------------------------------------------
\section{Methods}

This narrative review searched PubMed and Google Scholar for hangboard intervention studies to April 2026. Studies were included if they used a hangboard or no-hang device as the primary intervention, measured at least one finger-strength or endurance outcome, and had peer-reviewed publication status. Foundational biomechanical and physiological studies were included to contextualise training rationale. Practitioner claims from widely-used coaching resources (\textit{Training for Climbing}, \textit{Rock Climber's Training Manual}, Lattice Training, Dave MacLeod, Crimpd, \texttt{r/climbharder}) are reported and classified by evidence tier.

\textbf{Claim status labels.} Each claim carries one of four labels: \textbf{EVIDENCE\_BACKED} (peer-reviewed intervention evidence with transparent design); \textbf{BIOPLAUSIBLE} (mechanistically plausible, no direct controlled evidence); \textbf{FOLKLORE\_VERIFIED} (community-held belief with prevalence verified from $\geq$3 independent climbing-community sources); or \textbf{TO\_VERIFY} (provenance unclear or insufficient source verification; must not be treated as settled guidance).

\textbf{Evidence Strength Score (ESS).} Each actionable recommendation carries an ESS from 0 (unreliable source) to 10 (well-powered climbing RCT with functional outcomes). ESS reflects study design, sample relevance to climbers, risk of bias, and directness of the climbing application. See Appendix~\ref{app:toverify} for a full TO\_VERIFY tracker.

\textbf{Synthesis process.} Initial multi-domain literature queries used eight free-tier frontier LLMs to maximise recall across specialist topics (supplements, injury, recovery, weekly programming). Synthesis, claim classification, and ESS assignment were performed by the author using Claude Sonnet~4.6 as orchestration tool, with all claims reconciled against primary sources. Single-source claims carry an explicit footnote. Named model brands appear only in this methods paragraph and not in manuscript body text.
